\chapter{Conclusões}

O desenvolvimento da Proposta D exigiu um trabalho de integração significativo entre a lógica programável e o processador . Por limitações de tempo, não foi possível implementar todas as funcionalidades planeadas, como a interface DAC para emitir o sinal modulado e a validação final em placa. Ainda assim, o sistema desenvolvido até este ponto mostra que a arquitetura definida funciona como previsto.

O trabalho permitiu consolidar competências práticas de digital design entre PL e PS: desde a criação dos geradores e moduladores de sinal na FPGA até ao controlo de registos e parâmetros via barramento AXI no ambiente PYNQ . Isto cobre o objetivo principal da unidade curricular.
